% !TEX root = ./manual.tex
%%
% configurations
% @author Thomas Lehmann
%
\section{Konfiguration des Templates}\label{sec:configurations}
Das Template muss für die eigene Arbeit angepasst werden. Alle Einstellungen sind in der Datei \texttt{configuration/configuration.tex} vorzunehmen. Die Einstellungsmöglichkeiten und ihre Auswirkungen sind im folgenden aufgelistet.

\paragraph{Betreuer}\index{Einstellungen!firstSupervisor}\index{Einstellungen!secondSupervisor} Der Betreuer (Erstgutachter) und der Zweitgutachter werden über den Parameter der Kommandos \texttt{\textbackslash firstSupervisor\{\}} und entsprechend mit dem Kommando \texttt{\textbackslash secondSupervisor\{\}} eingestellt. Hier sind die beiden Default-Namen gegen die Namen der Gutachter auszutauschen. Ist kein Zweitgutachter beteiligt, dann muss der Parameter gelöscht werden.

\paragraph{Titel der Arbeit}\index{Einstellungen!termpaperTitle} Der eigentliche Titel der Arbeit wird im Parameter des Kommandos \texttt{\textbackslash termpaperTitle\{\}} angegeben. Ist der Titel länger und bildet drei Zeilen, so muss die Formatierung der Titelseite angepasst werden.

\paragraph{Keywords}\index{Einstellungen!Keywords} Für die Arbeit sind Schlüsselwörter in Englisch über Parameter des Kommandos \texttt{\textbackslash keywordsEN\{\}} und in Deutsch über das Kommando \texttt{\textbackslash keywordsDE\{\}} anzugeben. Die Begriffe werden unter dem Abstract eingefügt.

\paragraph{Autor}\index{Einstellungen!thesisAuthor} Der Autor wird im Parameter des Kommandos \texttt{\textbackslash termPaperAuthor\{\}} angegeben. Hier ist entsprechend der eigene Name einzusetzen.

\paragraph{Abgabedatum}\index{Einstellungen!submissiondate} Hier ist das Datum der Abgabe, so wie es im Dokument stehen soll, als Text im Parameter des Kommandos \texttt{\textbackslash submissionDate\{\}} anzugeben.

\paragraph{Non-Disclosure Agreement/Geheimhaltung}\index{Einstellungen!NDA} Ist die Arbeit Geheim, d.~h. es besteht ein Geheimhaltungsvertrag (NDA), so ist das Kommando \texttt{\textbackslash NDA\{yes\}} auszukommentieren. Ein entsprechender Vermerk erscheint dann auf dem Titelblatt. Ist die Arbeit nicht geheim, so ist \texttt{\textbackslash NDA\{no\}} auszukommentieren.

\paragraph{Cover-Page} Das Deckblatt ist über das Kommando
\texttt{\textbackslash Cover\{\}} mit entsprechendem Parameter
auszuwählen. Als Parameter sind zulässig:

\begin{center}
  \begin {tabular}{ll}
    Parameter & Cover Page \\
    \hline
    \texttt{Std2018} & Nachfolger des ''grünen Balkens'' \\
    \texttt{CD2017} & Corporate Design HAW Stand 2017 \\
    \texttt{CD2017NoLogo} & Corporate Design ohne Logos \\
  \end{tabular}
\end{center}

\paragraph{Art der Arbeit} Die entsprechende Arte der Semesterarbeit ist auszuwählen, indem das Kommando
\texttt{\textbackslash seminarKind\{\}} mit entsprechendem Parameter gewählt wird. 

\paragraph{Studiengang} Der Studiengang ist über das Kommando
\texttt{\textbackslash studycourse\{\}} mit entsprechendem Parameter
auszuwählen. Als Parameter sind zulässig:

\begin{center}
  \begin {tabular}{ll}
    Parameter & Studiengang \\
    \hline
    \texttt{INF\_ITS} & Bachelor Informatik Technischer Systeme \\
    %\texttt{TI} & Bachelor Technische Informatik \\
    \texttt{INF\_AI} & Bachelor Angewandte Informatik \\
    \texttt{INF\_WI} & Bachelor Wirtschaftsinformatik \\
    \texttt{INF\_ECS} & Bachelor European Computer Science \\
    \texttt{EMI\_EI} & Bachelor Elektrotechnik und Informationstechnik \\
    \texttt{EMI\_RE} & Bachelor Regenerative Energiesysteme und Energiemanagement \\
    \texttt{EMI\_IE} & Bachelor Information Engieering \\
    \texttt{BMT} & Bachelor Mechatronik \\
    \texttt{INF\_MaI} & Master Informatik \\
    \texttt{EMI\_MaA} & Master Automatisierung \\
    \texttt{EMI\_MaICE} & Master Information and Communnications Engieering \\
    \texttt{EMI\_MaMS} & Master Microelectronic Systems \\
  \end{tabular}
\end{center}

\section{Konfiguration im \LaTeX -Code}\label{sec:codechanges}
Für einige Anpassungen müssen direkt eine Änderungen im \LaTeX -Code vorgenommen werden. Soll eine Liste von Listings eingefügt werden, so ist die Zeile in der Datei \texttt{termpaper.tex} auszukommentieren und das Paket \texttt{listings} in der Paketlist hinzuzufügen.
